\section{Tecnologie}

Questa sezione descrive le principali tecnologie adottate nel progetto e le motivazioni che hanno guidato la loro selezione. L'obiettivo non e' fornire un semplice elenco di strumenti, ma mostrare come ogni scelta sia coerente con i requisiti del sistema e con l'architettura descritta nella sezione precedente.

\subsection{Stack frontend}

Il frontend di KompozeR e' stato concepito come una Single Page Application moderna, orientata a interazioni rapide e aggiornamento dinamico dello stato. In questa sottosezione e' opportuno descrivere le tecnologie utilizzate per:

\begin{itemize}
	\item rendering dell'interfaccia;
	\item gestione dello stato client;
	\item routing tra pagine e viste;
	\item integrazione con API REST e canali real-time.
\end{itemize}

Per ciascuna tecnologia selezionata conviene motivare la scelta in termini di produttivita', chiarezza architetturale, ecosistema e aderenza ai requisiti dell'applicazione.

\subsection{Stack backend}

Il backend e' organizzato come insieme di servizi web separati. Qui e' utile descrivere:

\begin{itemize}
	\item runtime applicativo;
	\item framework HTTP utilizzato per l'esposizione delle API;
	\item linguaggio e supporto alla tipizzazione;
	\item eventuale struttura interna comune dei microservizi.
\end{itemize}

In questa parte va spiegato perche' una soluzione basata su TypeScript e servizi Express risulta adeguata per un sistema web modulare e facilmente evolvibile.

\subsection{Persistenza dei dati}

La persistenza e' organizzata secondo il principio di database per servizio. Questa sottosezione dovrebbe spiegare:

\begin{itemize}
	\item la tecnologia di database scelta;
	\item le motivazioni di flessibilita' e velocita' di sviluppo;
	\item la separazione dei dati tra i diversi microservizi;
	\item l'allineamento tra modello logico del dominio e struttura dei documenti persistiti.
\end{itemize}

\subsection{Comunicazione real-time}

Alcune funzionalita' richiedono uno scambio di informazioni piu' immediato rispetto al classico modello request-response. In questa sottosezione e' utile presentare la tecnologia usata per la comunicazione real-time e spiegare in quali punti del sistema viene impiegata, ad esempio per notifiche utente e chatbot.

Conviene mantenere qui il focus sul valore applicativo e sull'esperienza utente, senza entrare nel dettaglio dei meccanismi distribuiti avanzati che saranno trattati nel report DS.

\subsection{Containerizzazione e strumenti di esecuzione}

Per il deployment locale e per l'avvio coordinato dei servizi e' utile descrivere gli strumenti utilizzati per containerizzazione e orchestrazione leggera. In particolare si possono motivare:

\begin{itemize}
	\item uso di Docker per isolare i componenti;
	\item uso di \texttt{docker-compose} per l'avvio dell'intero sistema;
	\item vantaggi in termini di replicabilita' dell'ambiente e semplicità di setup.
\end{itemize}

\subsection{Tecnologie per test e sviluppo}

Se previste o gia' utilizzate, questa sottosezione puo' descrivere le tecnologie di supporto allo sviluppo e alla qualita' del codice, ad esempio:

\begin{itemize}
	\item strumenti di test unitario e di integrazione;
	\item linters e formatter;
	\item strumenti per debugging e ispezione delle API.
\end{itemize}

\subsection{Valutazione complessiva dello stack}

In chiusura, conviene sintetizzare perche' lo stack scelto sia adeguato al progetto: buona separazione delle responsabilita', facilita' di sviluppo incrementale, coerenza con una web application moderna e supporto alla futura evoluzione del sistema.